\documentclass[11pt,a4paper]{article}

\usepackage[margin=1in]{geometry}
\usepackage{amsmath,amssymb,amsthm}
\usepackage{graphicx}
\usepackage{booktabs}
\usepackage{microtype}
\usepackage{xcolor}
\usepackage[hidelinks,colorlinks=true,linkcolor=blue!50!black,citecolor=blue!40!black,urlcolor=blue!50!black]{hyperref}
\usepackage[numbers,sort&compress]{natbib}
\usepackage{fancyhdr}
\pagestyle{fancy}
\fancyhf{}
\fancyhead[L]{\small\itshape Expansions of the Icosian Triad}
\fancyhead[R]{\small\thepage}
\fancyfoot[C]{\footnotesize\itshape Pre-peer-review open research preprint --- not independently validated.}
\renewcommand{\headrulewidth}{0.4pt}
\renewcommand{\footrulewidth}{0.4pt}

\newtheorem{theorem}{Theorem}
\newtheorem{proposition}[theorem]{Proposition}
\theoremstyle{definition}
\newtheorem{definition}[theorem]{Definition}
\theoremstyle{remark}
\newtheorem{remark}[theorem]{Remark}
\newtheorem{observation}[theorem]{Observation}

\newcommand{\Z}{\mathbb{Z}}
\newcommand{\Q}{\mathbb{Q}}
\newcommand{\R}{\mathbb{R}}
\newcommand{\C}{\mathbb{C}}
\newcommand{\II}{\mathcal{I}}
\newcommand{\GG}{\mathcal{G}}
\newcommand{\CC}{\mathcal{C}}
\newcommand{\NH}{\mathrm{N}_{\mathrm{H}}}
\newcommand{\Zphi}{\Z[\varphi]}
\newcommand{\WHfour}{W(H_4)}

\title{\textbf{Expansions of the Icosian Triad}\\[2pt]
\large Where $(\II, \GG, \CC)$ appears in cosmology, biology, cascade dynamics,\\[1pt]
and the existence-closure programme\\[6pt]
\small (Open research preprint --- pre-peer-review draft)}

\author{%
  Lee Smart\\[4pt]
  \textit{Institute of Vibrational Field Dynamics}\\[2pt]
  \texttt{contact@vibrationalfielddynamics.org}\\[2pt]
  \texttt{@vfd\_org}%
}
\date{May 2026}

\begin{document}
\maketitle
\thispagestyle{empty}

\begin{abstract}
\noindent
The companion paper~\cite{icosianTriadCore} (henceforth ``Paper 1'') derived
the icosian triad $(\II, \GG, \CC)$ from the geometry of the
$600$-cell vertex set $V_{600}$. The present paper catalogues the
\emph{expansions} of that triad — locations across the Vibrational Field
Dynamics programme where the same structural objects already appear, often
in seemingly unrelated contexts.

\smallskip\noindent\textbf{Five expansion domains.}
\begin{enumerate}
\item \textbf{Cascade ladder.} $V_{600}$ is the $H_4$ rung of the
cascade \cite{cascadeDerivation,cascadeRefinementPapers}. The
$\sigma$-fixed eigenspace family (total dimension $94$) and the
$\sigma$-paired eigenspaces ($E_1, E_2$ on one side, $E_6, E_8$ on the
other, with $4 + 9 = 13$ dimensions per side) are numerically adjacent
to the cascade's finite-level refinement infrastructure
\cite{cascadeRefinementPapers,cascadeSelectionH4}, but a vertex-level
intertwiner between Paper~1's $\C[V_{600}]$ eigenspaces and the
cascade's auxiliary fibre $F^0$ has not been constructed.

\item \textbf{Cosmology.} The hypersphere-cosmology release
\cite{hypersphereCosmology} uses the $+6\varphi$ adjacency eigenspace
of $V_{600}$ (Paper~1's $E_1$, dimension $4$) as the substrate of the
$\Lambda$ dipole and the $H_0$ correction. Paper~1 provides the formal
definition of that eigenspace.

\item \textbf{Biology / consciousness.} Note G \cite{noteG} states a
consistency check between volatile-anaesthetic potency and a candidate
microtubule-lattice realisation of the $+6\varphi$ dipole eigenspace
(Spearman $\rho = +0.94$, $n=11$ volatile inhaled agents only). The
substrate map remains conditional; the H-MT-Bridge cortical EEG
prediction is the independence-promoting test. Paper 1 provides the
underlying spectral object.

\item \textbf{Spectral bridge.} The Laplacian $\lambda_L = 12$
eigenspace, which is Paper~1's adjacency $E_4$ (dimension $25$, $A_1$
eigenvalue $0$), carries the Schl\"afli $5$-coset spectral bridge of
\cite{schlafliBridge,spectralBridge24600}. Throughout this paper the
convention is $\lambda_L$ for Laplacian $L = 12 I - A_1$ and
$\lambda_{A_1}$ for adjacency.

\item \textbf{Existence programme.} The two involution conventions of
\cite{existenceClosure} --- $\tau_\mathrm{ico}$ with
$\dim \mathrm{Fix} = 94$ and $\tau_\mathrm{spec}$ with
$\dim \mathrm{Fix} = 116$ --- relate to Paper~1 as follows.
$\tau_\mathrm{ico}$ is identified with Paper~1's Galois $\sigma$-action;
its $94$-dimensional fixed locus is the union of the $\sigma$-fixed
$A_1$-eigenspaces. $\tau_\mathrm{spec}$ is the spectral involution
acting as $-1$ on the adjacency $+6\varphi$ eigenspace (Paper~1's
$E_1$, dimension $4$) and as identity on the remaining $116$ dimensions;
a geometric or vertex-level realisation of $\tau_\mathrm{spec}$
remains open.
\end{enumerate}

\smallskip\noindent\textbf{Status.}
This paper makes no new mathematical claims about $V_{600}$. It is a
positioning document. Cross-references are to existing public preprints and
releases; predictions outside Paper~1 are stated as cited rather than as
re-derived.
\end{abstract}

\bigskip

\noindent\fbox{%
\parbox{0.96\linewidth}{%
\textbf{Plain-language summary.}

\smallskip\noindent
Paper~1 of this pair developed a single structural object — the
\emph{icosian triad} on $V_{600}$. Once that object is in hand, several
parts of the broader programme that had been described with different
languages turn out to be referring to the \emph{same} mathematical
content.

\smallskip\noindent
This paper points out those cross-identifications. It does not extend the
mathematics; it shows where the mathematics already lives in published
work, sometimes under different names.

\smallskip\noindent
The five recurrences listed in the abstract are: the cascade ladder,
cosmology, biology, the $24$--$600$ spectral bridge, and the existence
programme's involution conventions. Each is a separate paper or
preprint in the public record. Each uses, implicitly or explicitly,
the structure that Paper~1 derives in full.
}}

\bigskip

\subsection*{Expansion status table}

The five expansions below relate to Paper~1 with different status. The
following compact table records, for each:
the cited source, the type of mathematical / empirical content, whether
the cross-identification is a theorem, candidate substrate map, or
numerical adjacency, and the principal open gap.

\begin{center}\footnotesize
\renewcommand{\arraystretch}{1.25}
\begin{tabular}{|p{0.13\textwidth}|p{0.18\textwidth}|p{0.18\textwidth}|p{0.18\textwidth}|p{0.18\textwidth}|}
\hline
\textbf{Expansion} & \textbf{Source} & \textbf{Type} & \textbf{Status} & \textbf{Principal open gap} \\
\hline
Cascade ladder
  & \cite{cascadeRefinementPapers,cascadeSelectionH4}
  & Cited theorems (28-dim $\WHfour$-fixed $F^0$, block Laplacian)
  & Spectral adjacency to Paper~1's eigenspaces, no intertwiner
  & Explicit intertwiner $F^0 \leftrightarrow \C[V_{600}]$ \\
\hline
Cosmology
  & \cite{hypersphereCosmology}
  & Empirical, 7 observables $< 0.5\%$ Planck 2018
  & Candidate substrate map; uses $E_1$ dipole numerically
  & Operator-level identification of cosmology $F$ with $\CC_\varphi|_{E_1}$ \\
\hline
Biology
  & \cite{noteG}
  & Spearman $\rho = +0.94$ consistency check (volatile, $n=11$)
  & Candidate microtubule substrate map; H-MT-Bridge conditional
  & Cortical-EEG $R_{\mathrm{phase}}$ pre-registered test \\
\hline
Spectral bridge
  & \cite{schlafliBridge,spectralBridge24600,bridgeSynthesis}
  & Cited theorems (exact $\Q(\sqrt 5)$ arithmetic)
  & $\lambda_L = 12$ lift into Paper~1's $E_4$
  & $A_5$ vs $\WHfour$ equivariance distinction --- settled in this paper \\
\hline
Existence programme
  & \cite{existenceClosure,lifeAsClosure}
  & Cited definitions ($\tau_{\mathrm{ico}}$, $\tau_{\mathrm{spec}}$)
  & $\tau_{\mathrm{ico}}$ = Paper~1's $\sigma$; $\tau_{\mathrm{spec}}$ spectrally defined as $-1$ on $E_1$
  & Geometric / vertex-level realisation of $\tau_{\mathrm{spec}}$ \\
\hline
\end{tabular}
\end{center}

This table is the load-bearing scope artefact for the cross-identifications
below. Each subsequent section should be read against its row.

\bigskip

\section{Introduction}\label{sec:intro}

The icosian triad $(\II, \GG, \CC)$ defined in Paper~1 is the structural
object obtained by combining three simultaneous facts about the unit
icosian group $V_{600}$ inside the quaternion algebra over $\Q(\sqrt 5)$:

\begin{itemize}
  \item a $9$-class symmetric Bose--Mesner association scheme with
    $A_1$-spectrum in $\Zphi$ and multiplicity sequence
    $(1,4,9,16,25,36,9,16,4)$;
  \item a closure operator $\CC_\varphi = L_{V_{600}} + \varphi^{-2}\!\cdot\! I$
    acting diagonally on the $9$ eigenspaces;
  \item a quaternion norm map $\GG = \NH : \II \to \Zphi$ surjective onto
    totally positive $\Zphi$-primes, with $\sigma$-equivariance
    $\sigma \NH = \NH \widehat\sigma$.
\end{itemize}

The total dimension $\dim \mathrm{Fix}(\sigma) = 94$ on the $A_1$-eigenspace
decomposition is the load-bearing invariant.

The Vibrational Field Dynamics programme contains five public preprints (or
release bundles) in which this triad already appears in disguise. The present
paper catalogues these appearances. The mathematical content of each
appearance is in the cited source; here we connect them via Paper~1.

\bigskip

\section{Expansion 1: the cascade ladder}\label{sec:cascade}

\subsection*{Statement}

The cascade ladder \cite{cascadeDerivation,cascadeRefinementPapers}
places $V_{600}$ at the $H_4$ rung, with downstream rungs at the
$40$-vertex root system (cf.\ \cite{cascadeRefinementPapers}) and below.
The cascade-selection-h4 paper \cite{cascadeSelectionH4} requires a
$\WHfour$-fixed subspace $(F^0)^{\WHfour}$ of dimension $28$ on the
fibre of each level of the finite-level refinement, and witnesses the
$(\mathrm{O}1)$ closure obligation via the block Laplacian
$L_N^\mathrm{block} = L_N^\mathrm{cl} \otimes \mathrm{id}_{F^0}$.

\subsection*{Connection}

Paper~1's $9$-class scheme contains the $\WHfour$-representation structure
of $V_{600}$ at the $\C[V_{600}]$ level. The $28$-dimensional
$\WHfour$-fixed subspace of the cascade fibre $F^0$ in
\cite{cascadeSelectionH4} is a separate structural object from the
$\C[V_{600}]$ eigenspaces; an explicit intertwiner between the two
spaces is not yet constructed.

What can be cleanly stated: cascade-selection-h4's block-Laplacian
lowest non-zero eigenvalue at $N=0$,
$\lambda_+ = 9 - 3\sqrt 5 = 12 - 6\varphi$, equals the Laplacian
$L = 12 I - A_1$ eigenvalue corresponding to Paper~1's $A_1$-eigenvalue
$\lambda_1 = 6\varphi$ (eigenspace $E_1$, dimension $4$). The
$\sigma$-conjugate pair $(E_1, E_8)$ of Paper~1's spectrum is the
Galois partner that organises this lowest cascade band; explicit
identification of the fibre $F^0$ with subspaces of $\C[V_{600}]$ is
deferred to a follow-on note.

\subsection*{Consequence}

Paper~1 provides the full $V_{600}$ spectral basis that organises the
cascade's $H_4$-rung structure at the polytope level. The explicit
intertwiner between cascade fibre $F^0$ and $\C[V_{600}]$ remains an
open structural question.

\bigskip

\section{Expansion 2: cosmology}\label{sec:cosmology}

\subsection*{Statement}

The hypersphere-cosmology release \cite{hypersphereCosmology} derives
the cosmological constant $\Lambda$ as a $\sigma$-paired dipole
correction to the $E_8$-substrate quadratic form, with the Planck-side
numerics $\Lambda$ within $0.078\%$, $H_0$ within $0.45\%$, and
$\Omega_\Lambda$ within $0.083\%$ of Planck 2018 inferred values,
together with four further cosmological observables to $< 0.5\%$ (seven
total). The dipole construction uses the adjacency eigenvalue $+6\varphi$
sector on $V_{600}$ (Paper~1's $E_1$, dimension $4$) and the
complementary Laplacian value $12 - 6\varphi$.

\subsection*{Connection (candidate substrate map)}

The hypersphere-cosmology dipole sector is, numerically, the
adjacency-$+6\varphi$ eigenspace of $V_{600}$, equivalently Paper~1's
$E_1$ (dimension $4$). The Galois partner $E_8$ (adjacency $6 - 6\varphi$,
dimension $4$) does not enter the cosmology paper's $\Lambda$ formula
directly; it is recorded by the Galois pairing as a separate object.
The full identification of the cosmology source's operator
$F = \alpha R + \beta E - \gamma Q$ with a restriction of Paper~1's
$\CC_\varphi$ is conditional on the cascade-substrate map and is not
established at the operator level in either paper.

\subsection*{Consequence}

Paper~1 supplies the spectral basis ($E_1$ at adjacency $+6\varphi$) in
which the cosmology source's dipole construction lives. The cosmology
numerical success ($94/94$ verification tests) is consistent with the
icosian triad as candidate substrate; a full operator-level
identification would convert the consistency into a structural prediction
and is open.

\bigskip

\section{Expansion 3: biology / consciousness}\label{sec:biology}

\subsection*{Statement}

Note G \cite{noteG} reports that the $\sigma$-paired closure-dipole
class on $V_{600}$, when interpreted as a candidate microtubule-lattice
dipole substrate, gives a consistency check between
literature-DFT tubulin dipoles and clinical anaesthetic potencies:
Spearman $\rho = +0.94$ on $n = 11$ volatile inhaled anaesthetics (the
all-agent correlation including IV agents is much weaker,
$\rho \approx +0.5$). Note G frames this as a consistency check, not
independent causal evidence. Note G also states the H-MT-Bridge
conditional proposition: the cortical EEG $R_{\mathrm{phase}}$ ratio
(sevoflurane / propofol) is at least $1.5$ under equipotent dosing.

\subsection*{Connection (candidate substrate map)}

The dipole class of Note G is, in Paper~1's notation, the same
adjacency-$+6\varphi$ eigenspace $E_1$ (dimension $4$) that appears in
the cosmology source. The substrate map from Paper~1's $E_1$ to the
tubulin $\alpha$/$\beta$-dimer dipole array is a conditional candidate
mapping; Note G explicitly does not validate it.

The Note~G T15 result is a uniform $V_{600}$-leverage observation, not a
residue-intersection result. The biological closure-control-node
candidate residues identified in \cite{noteG} arise from downstream
symmetry breaking on the candidate substrate, not directly from the
$E_1$ basis.

\subsection*{Consequence}

Paper~1 supplies the spectral object ($E_1$ at adjacency $+6\varphi$)
that Note G uses as its candidate substrate. The $\rho = +0.94$
correlation is a consistency check on the candidate map (volatile-agent
subset). The H-MT-Bridge cortical-EEG prediction is the
independence-promoting test; until that test produces data, the substrate
identification stays conditional.

\bigskip

\section{Expansion 4: the 24-600 spectral bridge}\label{sec:24600}

\subsection*{Statement}

The three-paper $24$--$600$ spectral bridge
\cite{schlafliBridge,spectralBridge24600,bridgeSynthesis} establishes:
\begin{itemize}
  \item the $5$-coset decomposition $V_{600} = 2T \sqcup 2T\!\cdot\! g \sqcup 2T\!\cdot\! g^2 \sqcup 2T\!\cdot\! g^3 \sqcup 2T\!\cdot\! g^4$
    (right cosets; $2T$ is non-normal in $V_{600}$) into five $24$-cell cosets;
  \item the selective lift of each local $24$-cell Laplacian
    $\lambda = 12$ eigenspace (dimension $2$) into the global $V_{600}$
    Laplacian $\lambda = 12$ eigenspace, of dimension $25$, with exact
    $\mathbb{Q}$-rational characters;
  \item an explicit $A_5$ action on the cosets with character
    $5 \cdot Y_5 = 2 \cdot Y_5 \oplus 3 \cdot Y_5$.
\end{itemize}

Throughout this section the symbol $\lambda$ refers to the
Laplacian $L = 12\, I - A_1$, not to the adjacency $A_1$ used in
Paper~1's spectrum table.

\subsection*{Connection}

Paper~1 supplies the full $A_1$-eigenstructure of $V_{600}$.
Translating between conventions: the bridge's
Laplacian eigenvalue $\lambda_L = 12$ corresponds to Paper~1's
$A_1$-eigenvalue $\lambda_{A_1} = 0$, hence to eigenspace $E_4$
(dimension $25$). The bridge's global $25$-dimensional $\lambda_L = 12$
space is therefore exactly Paper~1's $E_4$. Separately, Paper~1's
$E_0$ (dimension $1$, the all-ones eigenspace at $A_1$-eigenvalue $12$,
equivalently Laplacian eigenvalue $0$) accounts for the global constant
row inside this picture.

The spectral bridge's selective lift corresponds to the
$A_5$-equivariant projection onto $E_4$ from each local
$24$-cell Laplacian's two-dimensional $\lambda_L = 12$ subspace.

\subsection*{Consequence}

The spectral bridge is the $E_4$-projection of a broader $9$-eigenspace
structure, with $E_0$ accounting separately for the constant row.
Paper~1 places this in the wider context. The other seven eigenspaces of
$V_{600}$ ($E_1, E_2, E_3, E_5, E_6, E_7, E_8$) carry the $\sigma$-paired
dipole class, the antipodal $\mathbb{Z}_2$-grading, and the
$\sigma$-paired $9$-dimensional content, none of which is visible from
the spectral bridge alone.

\bigskip

\section{Expansion 5: the existence programme involutions}\label{sec:existence}

\subsection*{Statement}

The existence-life-closure programme \cite{existenceClosure,lifeAsClosure}
uses two involutions on $V_{600}$:
\[
  \tau_\mathrm{ico} \;:\; \dim \mathrm{Fix} = 94, \qquad
  \tau_\mathrm{spec} \;:\; \dim \mathrm{Fix} = 116.
\]
Both commute with the closure operator $\CC_\varphi$ and both appear in
theorem statements (the former) and simulation demonstrations (the latter)
of the existence programme.

\subsection*{Connection}

$\tau_\mathrm{ico}$ is the Galois action $\sigma$ of Paper~1. Its fixed
locus, restricted to the $A_1$-eigenspace decomposition, is
\[
  \mathrm{Fix}(\sigma) \;=\; E_0 \oplus E_3 \oplus E_4 \oplus E_5 \oplus E_7
  \;=\; \dim 1 + 16 + 25 + 36 + 16 = 94,
\]
exactly matching the existence programme's load-bearing dimension.

\subsection*{Consequence}

Paper~1 derives the $94$-dimensional $\sigma$-fixed locus structurally. The
existence programme's $\tau_\mathrm{ico}$ convention is the same object.

\subsection*{Spectral status of $\tau_\mathrm{spec}$}

The second involution $\tau_\mathrm{spec}$ with
$\dim \mathrm{Fix} = 116$ is spectrally identified in
\cite{existenceClosure}: it acts as $-1$ on the adjacency $+6\varphi$
eigenspace (Paper~1's $E_1$, dimension $4$) and as identity on the
remaining $116$ dimensions. What remains open is a geometric /
vertex-level realisation of this involution; that is a clean
combinatorial follow-on.

\bigskip

\section{What this paper does not do}\label{sec:non-claims}

\begin{itemize}
  \item This paper does not extend Paper~1's mathematics. Quantitative
    claims used here are either cited from a prior source (cosmology,
    biology, cascade, bridge papers) or are direct arithmetic
    translations of Paper~1's spectrum into those sources' conventions.
    Cross-domain identifications are positioning claims, conditional on
    the substrate maps stated in each section.
  \item It does not address the Riemann Hypothesis or its generalisations.
  \item It does not propose new empirical predictions. The predictions
    referenced (e.g.\ Note G's H-MT-Bridge) remain at the status established
    in their source documents.
  \item It does not aim to be the canonical synthesis of the VFD
    programme; it is a localised positioning paper for the icosian triad's
    appearances in cosmology, biology, cascade dynamics, the
    $24$--$600$ spectral bridge, and the existence-closure programme.
\end{itemize}

\bigskip

\section{Open follow-on directions}\label{sec:open}

\begin{enumerate}
  \item Geometric identification of $\tau_\mathrm{spec}$ (existence
    programme) as an involution on $V_{600}$ with $\dim \mathrm{Fix} = 116$.

  \item Construction of explicit Brandt / Hecke operators at higher prime
    powers; decomposition of $\C[V_{600}]$ into Hecke-eigen-pieces.

  \item Refinement of Note G's H-MT-Bridge from a consistency check
    into an operator-level prediction tied to Paper~1's $E_1$
    eigenspace (adjacency $+6\varphi$) and its Galois partner $E_8$.

  \item Numerical study of zeros of $L(\Theta_\II, s) = \zeta_K(s) \zeta_K(s-1) \cdot C_2(s)$ (where $C_2(s)$ is the explicit local-$2$ factor of Paper~1, Theorem~18)
    on the critical lines $\mathrm{Re}(s) = 1/2, 3/2$ (subject to classical
    RH and GRH for $\chi_5$).

  \item Comparison of the icosian theta with other definite quaternion-order
    theta functions over real quadratic fields (Hurwitz over $\Q$, Lipschitz
    over $\Q$, octahedral order over $\Q(\sqrt 2)$, $\dots$) to characterise
    what is unique about the $\Q(\sqrt 5)$ / icosian case.
\end{enumerate}

\bigskip

\bibliographystyle{plainnat}
\begin{thebibliography}{99}

\bibitem{icosianTriadCore}
L.~Smart,
\emph{The icosian triad on $V_{600}$: a geometry-first derivation},
Companion to this paper, Institute of Vibrational Field Dynamics, 2026.

\bibitem{cascadeDerivation}
L.~Smart,
\emph{Cascade derivation: from $E_8$ to the polytope ladder},
Open research preprint, Institute of Vibrational Field Dynamics, 2026.

\bibitem{cascadeRefinementPapers}
L.~Smart,
\emph{cascade-refinement-papers (P3 + P4 + cascade-refinement-compat + cascade-mechanism + cascade-selection-h4)},
Five-paper bundle, Institute of Vibrational Field Dynamics, 2026.

\bibitem{cascadeSelectionH4}
L.~Smart,
\emph{cascade-selection-h4: a witnessed $(\mathrm{O}1)$ discharge on the $H_4$ branch},
In~\cite{cascadeRefinementPapers}.

\bibitem{hypersphereCosmology}
L.~Smart,
\emph{Hypersphere Cosmology v1.1.0},
Public release, github.com/vfd-org/hypersphere-cosmology, 2026.

\bibitem{noteG}
L.~Smart,
\emph{Note G: Tubulin dipoles and the $\sigma$-paired closure-dipole class},
Open research preprint, Institute of Vibrational Field Dynamics, 2026.

\bibitem{schlafliBridge}
L.~Smart,
\emph{The Schl\"afli decomposition of the 600-cell},
In~\cite{spectralBridge24600}.

\bibitem{spectralBridge24600}
L.~Smart,
\emph{the-24-600-spectral-bridge (three-paper bundle)},
Public release, github.com/vfd-org/the-24-600-spectral-bridge, 2026.

\bibitem{bridgeSynthesis}
L.~Smart,
\emph{From Schl\"afli decomposition to spectral bridge: a first closure-projection channel in $V_{600}$},
In~\cite{spectralBridge24600}.

\bibitem{existenceClosure}
L.~Smart,
\emph{Existence as Closure},
Open research preprint, Institute of Vibrational Field Dynamics, 2026.

\bibitem{lifeAsClosure}
L.~Smart,
\emph{Life as Closure},
Open research preprint, Institute of Vibrational Field Dynamics, 2026.

\end{thebibliography}

\end{document}
